\section{第四原则：递归优化原则}

前面三条原则处理边界、标准与分工。本节继续检查维持它们的更新过程：实践带来证据，证据推动标准修订，修订又改变下一轮行动。回路可以形成稳定状态，也可以随环境持续变化，具体行为取决于更新规则、延迟与外部输入。

把这个回路写成迭代的形式，它的行为由收敛条件决定，

\begin{equation}
x_{n+1}=g(x_{n}),\qquad \left|g'(x^{*})\right|<1\ \text{时迭代收敛到不动点}\ x^{*}
\end{equation}

这里先以一个标量参数$x_n$说明评价更新，$g$为连续可微的更新映射，$x^*$为不动点。$|g'(x^*)|<1$给出局部渐近收敛的充分条件；导数绝对值大于1时，不动点不稳定，等于1时还须进一步分析。多参数更新则要考察雅可比矩阵的特征值。价值与环境随实践改变，可能形成新的稳定状态，也可能持续变动，具体行为由更新规则和外部扰动决定。

\subsection{反向算法：主动引入摩擦}

\subhead{问题的位置}

前面的讨论给出的是价值流向和人的重新定位，这些是结构层面的判断。结构层面的判断有一个弱点：它正确，但不能直接指导一个人在星期一早上做什么。压力落在个体身上时，不会以结构分析的形式出现，而是以一连串具体选择的形式出现：用不用导航，写不写这段文字，回不回这条消息。这一节处理的就是这些选择，但处理方式不是给出一份清单。

先说清楚这一节不做什么。它不主张拒绝技术，也不主张人为制造困难。一个人的日常决定散布在无数细小的场景里，任何一份统一的清单都会在第一次遇到例外时失效。这里要处理的是一个更窄的问题：低黏滞系统会持续把人的行为推向某些固定方向，如果一个人想在这些方向上保留一部分自主，需要满足什么条件。这不是关于意志的问题，而是关于边界维护的问题，因此它可以用结构的方式来说。

\subhead{系统把人推向哪里}

要保留自主，先要看清推力指向哪里，因为反抗一个方向的前提是知道方向。低黏滞系统对人行为的塑造可以归纳为四个方向，而它们各自都是被优化出来的，不是被设计出来的。第一个方向是最短路径。只要存在一条更短的路径，系统就会把它找出来并设为默认，用户每一次手动绕开都在付出额外成本。第二个方向是即时反馈。反馈延迟越长，用户越容易流失，因此所有环节都在被压缩，包括那些本来需要时间的环节。

第三个方向是偏好闭环。系统从一个人的历史行为中推断他的偏好，然后按这个偏好供给内容，供给的结果又成为新的历史行为，闭环因此逐步收窄。这个过程不需要任何恶意，它只是让预测更准的必然结果，而更准的预测意味着更少的意外。第四个方向是无痛满足。凡是能提供的满足，系统都倾向于以最低成本提供，因为低成本意味着更高的使用频率，而高成本的部分被逐步剥离。四个方向的共同点是它们都提高短期指标，因此不需要任何人推动，它们自己就会发生。

这些推力不是外来的胁迫，它们不禁止任何事。一个人随时可以关掉推荐、可以选择手动路线、可以拒绝生成。问题在于每拒绝一次都要付出成本，而这些成本是持续的、累加的、没有外部奖励的。在一个推力持续存在而阻力没有回报的环境里，行为的分布会向推力方向偏移，这不是判断力的问题，是概率的问题，除非有某种机制能稳定地抵消它。

\subhead{保留非最优接口}

能稳定抵消它的机制，是把一部分非最优的选择固定下来，使它不再需要每次重新决定。这就是所谓保留非最优接口的意思。接口在这里指的是人与世界接触的那一层：怎么走一段路，怎么读一本书，怎么完成一件东西，怎么和一个人维持联系。把某些接口刻意保留在非最优状态，代价是明确的效率损失，收益是这些接口所承载的能力不会退化。

具体的形式因场景而异，但有几类被反复验证过。第一类是等待。主动接受某些延迟，不用缓冲把它填满，让间隙重新成为一个可以思考和退出的位置。第二类是身体训练。它是最不容易被代劳的接口，因为肌肉、心肺和协调性的改变必须由本人的重复动作产生，而且这个过程无法被加速，训练量的积累和能力的提高之间有硬性关系。第三类是线下关系。面对面接触携带的信息量和约束强度高于远程交流，而它的成本也更高，因此是第一批被替代的接口。

还有几类性质稍有不同的。手作与制作把人和材料直接接上，材料会抵抗，会出错，会留下不可撤销的痕迹，这些性质在数字工具里被全部移除。难读的文本要求持续的注意力，不能跳读，不能摘要，读者必须自己完成理解过程。长期项目在几个月到几年的尺度上没有任何阶段性成就，它的回报集中在完成时刻，因此和即时反馈系统正好相反。无即时反馈的承诺要求一个人在没有回报的期间继续兑现。偶遇与随机性把非最优的结果保留下来，使偏好闭环无法完全闭合。

这些活动的价值在于它们仍能形成所需要的能力。若直接练习比完全代办更能维护判断，就可以把练习安排成稳定的默认；若工具提供更好的反馈，也应采用。相对即时产出的效率损失，需要与长期能力收益一起比较。

\subhead{为什么不是意志力问题}

把上面的做法理解为自律，会引出两个错误的推论。第一个是把它当成道德要求：能坚持的人更有品格，坚持不了的人是自己不行。第二个是把它当成个人偏好：愿意怎样生活是各人的选择，不需要理由。两个推论都错在同一个地方，它们把这件事放在个人层面处理，而它的问题出在环境的反馈结构上。

低黏滞环境不是一个人可以选择退出的外部条件，它已经嵌入了工作、社交、支付、出行和求医的每一个环节。在这种条件下，要求一个人靠意志力持续逆着环境走，等于要求他一生都在消耗同样的资源，而且这个消耗不会带来任何外部承认。这个要求是不现实的，不是因为它太难，而是因为它没有稳定的支撑结构。真正有效的做法不是在意志层面加码，而是在结构层面减少需要意志介入的次数：把某些接口从选项中移除，让它们不进入每次决策的比较。

这个区分的实践含义是具体的。一个把手机放在另一个房间的人，并不是在考验自己的意志，他是在改变自己需要做出的决定数量。一套固定的作息安排之所以比每天决定几点睡更有效，也是因为它把决定变成默认。这些做法在表面上像是自律技巧，实质上是把一部分选择权从当下的自己手里拿走，交给一个事先设定好的安排，从而抵抗环境对行为分布的持续拉扯。

\subhead{代价的方向}

需要把这一节的限度说清楚，否则它会变成一份生活方式建议，而生活方式建议在结构上站不住。这些做法的代价是真实且不对称的。保留非最优接口的人会在具体的场合输给采用最优路径的人：更慢，更累，产出更少，而且在大多数评价体系里得不到任何补偿。这个代价是持续的，而收益是长期的、分散的、难以指认的，因此在个人层面很难形成稳定的正反馈。

这意味着这一节的建议不构成一种普遍可行的方案。对能够承担这个代价的人，它是有效的；对时间、精力或经济条件不允许的人，它是奢侈的，把它普遍化会变成对处境不利者的又一次指责。能承担的范围内做多少，取决于具体情况，而更重要的是认识到一件事：个体层面的调整只能减缓偏移的速度，不能改变偏移的方向。方向由评价函数决定，而评价函数的改变不是个体能完成的动作。把这一层说清楚，是为了避免把结构问题伪装成个人修养问题，也是为了指出下一步该往哪里走：从个体转到组织，从个人选择转到那些决定评价函数的地方。

\subsection{高黏滞组织：信任与验证的重建}

\subhead{从个体到组织}

把个体的调整做到限度之后，会遇到一个不由个人决定的东西：方向由评价函数决定，而评价函数不由个人设定。评价函数存在于组织中，因此改变它需要在组织层面完成。这一节处理的就是这类组织，但切入的角度不是管理，而是一个具体的功能问题：在一个生成成本趋近于零的环境里，什么样的组织还在生产别的东西。

要看清这个问题，可以先看一个反差。同样面对生成式工具，两类组织的反应完全不同。一类组织立刻把它接入生产流程，交付速度上升，人力需求下降，报表改善；另一类组织允许使用，但同时增加了验证环节、收紧了发布权限、延长了关键判断的复核周期，短期内效率反而下降。按通常的效率标准，第一类组织是赢家。但如果把时间尺度拉长，并且把事故率、召回率、返工率、诉讼成本和人员流失一起计入，两类的排序可能会反过来。这个反转会不会发生，取决于该组织的产出在多大程度上依赖判断的正确性，以及错误一旦发生的代价有多大。

\subhead{凭据生成器}

那些排序可能反转的组织，有一个共同的功能，可以称之为凭据生成。凭据在这里指的是一个可以被第三方依赖的判断：这份检测报告可信，这个药可以上市，这座桥能承重，这份审计意见可以拿去贷款，这个判决可以执行，这篇论文的结论可以被引用。凭据的价值不在于它有多新颖，而在于它背后的判断已经经过了某种检验，而且这个检验的过程可以被追溯和复核。

大学、实验室、医院、工程机构、法院、审计机构、认证机构、事故调查委员会，都属于凭据生成器。它们看起来效率低下：审稿要几个月，临床试验要几年，工程验收要停工期，判决要经过多级审理，事故调查要反复取证。这些延迟经常被批评为官僚主义，其中也确实有一部分是。但把延迟整体删除之后，这些机构就不再生产凭据，它们只生产意见，而意见不能作为凭据被依赖，因为它没有承载责任的构造。

理解这个功能之后，就能看清这些机构的慢在哪里是有意义的。有意义的慢集中在判断的形成环节：证据如何收集，反方意见如何进入，结论如何被复核，谁有权签署。无意义的慢集中在行政流转环节，这些环节的删除通常不损害凭据质量。可信组织的价值，不在于生成更快，而在于让判断有地方对账。对账的意思是：当一个判断被质疑时，有能力回到形成这个判断的过程，并检查每一步是否成立。

\subhead{四类慢变量}

凭据生成依赖于几类慢变量，它们各自对应一种不能压缩的过程。第一类是验证周期。一个结论要被确认为可靠，需要经过独立的重复，而重复需要时间。药物的三期临床、材料的长期老化测试、算法的对抗测试、结构的安全系数验证，都属于这一类。验证周期可以被组织得更好，但不能被跳过，因为跳过之后得到的不是更快的结论，而是一个未被验证的结论。

第二类是信任积累。一个机构的可信度由它过去的行为记录决定，而这个记录只能通过时间生成。一家医院需要多年的诊疗记录才能建立转诊信誉，一个认证机构需要经手足够多的案例才能被市场接受，一个实验室需要若干篇经得起复查的论文才能获得同行信任。信任的积累慢而摧毁快，这个不对称使它对机构的激励结构提出了很高的要求：任何一次为了短期收益而放松的验证，都可能把多年积累一次性抵消。

第三类是责任链的清晰度。判断由谁做出、依据什么、在什么权限范围内、出了错由谁承担，这些问题必须在判断形成之前就有答案，事后补是不行的。事故调查中最常见的问题不是找不到责任人，而是责任链在制度设计阶段就不存在，导致事后无法归属。第四类是慢变量本身的保护机制。承上所述，评价体系天然偏好可量化的短期指标，因此那些需要长期投入的环节会被持续挤压，除非有明确的机制保护它们：终身教职、独立预算、法定的复核程序、行业自律规范，这些安排的历史来源各不相同，但功能相同：让某些环节不必为当期的产出指标负责。

\subhead{为什么市场不会自动提供}

如果凭据生成这么重要，一个自然的问题是为什么市场不自动产生足够的供给。这个问题有几个层次的答案。第一层是时间错配。凭据的价值在被需要时才显现，而事故是低频事件，因此在两次事故之间的长周期里，投入凭据生成的组织看起来只是在浪费钱。没有经历过事故的决策者很难为它付账，而经历过事故的决策者往往已经不在位置上了。

第二层是收益的不可独占性。一个机构投入成本建立高标准的验证流程，它生产出的可信度有一部分会外溢给整个行业：客户会把整个行业的标准提高，而不只是信任这一家。这种正外部性使个体投入低于社会最优水平，这是认证、标准、审计这类活动长期需要公共或行业层面强制的原因。第三层是竞争压力下的向下竞次。当同一市场中有一方降低了验证标准并把节省的成本让给客户，其余各方在短期内有强烈的动机跟随，而后果的显现要等到下一次事故。这个结构在金融、航空、建筑、食品、医药领域都反复出现过，每一次的应对方式都是事后监管，而监管的时滞本身就是代价的一部分。

这些机制说明，凭据生成不能完全靠市场自发提供，但它也不能完全靠行政命令提供，因为凭据的可信度恰恰来自生成机构的独立性。这个张力没有一般解，只能在具体领域里通过制度安排来平衡，而平衡点在不同领域差别很大。

\subhead{组织内部的改变}

回到最初的问题：一个组织要维持凭据生成能力，需要在自己内部做什么。最容易想到也最容易被做错的，是设置验证环节。单纯增加一道审批不生产凭据，它只生产签字，因为审批者如果没有足够的信息和动机去实质检验，他的签字只是在分摊责任，不是在确认判断。有效的验证需要三个条件同时满足：验证者能看到原始材料而不是加工后的摘要，验证者的评价不受被验证方的直接控制，以及验证者做出否定判断时不会受到惩罚。

除了验证，还有两处需要调整。一是把慢变量的保护写进考核。如果所有指标都是季度产出，那么长期投入必然被牺牲，这不是执行问题，是指标结构问题，唯一的解法是让一部分评价周期长于任期。二是保留可以被追溯的记录。生成式工具让产出变得容易，也让产出的来源变得难以分辨，因此一份记录了谁在什么依据下做了什么判断的档案，其价值在上升。这类记录在平时被视作负担，在事故发生时是唯一能重建过程的东西。

需要承认这套安排的成本。一个维持凭据生成能力的组织，在生成速度上会持续落后于不维持的组织，而这个落后是可见的、可量化的、会在竞争中被反复计入的；它换来的东西是不可见的，因为它体现在没有发生的事故里。这个不对称是这类组织长期面临的困境，也是它们需要外部制度支持的原因。个体在这一层能做的有限，但组织层面的选择决定了评价函数，而评价函数决定了前面那些个体做法能不能长期维持。如果组织不保护慢变量，个体的保留非最优接口就只能在私人时间里进行，而私人时间也是被同一套评价函数压缩的。

到这里，个体和组织各自的动作都已经交代清楚，但它们都建立在一个尚未检查的假设上：人被当作产出能力来管理，因此保留位置的问题被理解成保留产出能力的问题。这个假设看起来很自然，实际上它有一个可以追溯的来历，而把这个来历看清楚，前面两节的做法才有可能被放在正确的位置上。

\subsection{负熵体验者：功能性与体验性的正交}

\subhead{一个没有被检验的前提}

如果人承担的是定义地形的角色，那么人自身应当被理解成什么，就成为一个必须重新回答的问题。把人当作产出工具的旧框架，在这个位置上已经不够用了；这个判断需要一个正面的替代，而在给出替代之前，先要看清楚被替代的是什么。关于人工智能与人的讨论，绝大多数都建立在一个没有明说的前提上：人的价值等于人的产出。这个前提如此普遍，以至于争论的双方都在用它。乐观的一方说机器会把人从重复劳动中解放出来，让人去做更有创造性的产出；忧虑的一方说机器会取代人的产出能力，使人失去位置。两种说法方向相反，前提相同。

如果这个前提是对的，那么当产出能力被普遍替代，人的位置问题就只剩下两个可能：找到一个机器还做不了的产出，或者接受自己不再必要。这两种前景都不好受，也都不准确，因为它们都接受了一个未经检验的前提。问题出在前提本身。人是产出工具，这个理解不是从人的本性中读出来的，它是工业化度量体系的一个副产品，而副产品是可以被追溯的。把这个前提的来路看清楚，它的约束力就会松动。

\subhead{度量体系的副产品}

国民核算体系在二十世纪中叶成型，它需要一个共同的度量单位来比较不同的活动，最终选择了货币化的产出。这个选择带来了巨大的便利：不同质的东西可以被加总，增长可以被比较，政策效果可以被评估。它同时带来一个后果：凡是不进入产出计量的活动，在核算中不存在。这个后果不是被有意识设计的，但影响深远。家务劳动、照护、社区互助、非正式的知识传递，长期不在核算范围内，因此长期不被当作经济事实。这些活动的大部分承担者是女性，这个分布使问题被进一步忽视。直到今天，不同国家对无偿照护的估算差异仍然很大，而这些估算的范围本身就说明，我们对什么算工作的定义是构造出来的。

人被理解成产出工具，正是这个过程的结果。当产出成为唯一被系统记录的东西，人就被窄化成了产出能力，而人的其他部分（他的时间、他的关系、他的有限性、他的经历）因为不进入记录，慢慢退出了公共讨论的视野。这个窄化在产出稀缺的时代不明显，因为那时候讨论人的位置就是在讨论生产能力，两者恰好重合，把两者等同起来几乎不产生误差。当生成能力变得充足，重合消失了，前提的缺陷才暴露出来。换句话说，这个前提一直是错的，只是过去一直没有被检验的机会。

\subhead{不能被折算的部分}

一个人除了处理信息、搬动物质、完成任务，还在经历别的事情。这些事情的共同特征是它们不能被折算成产出单位，因此在当前的核算体系里不留下痕迹，而它们并非无关紧要的剩余。痛苦是一个。它不产出任何东西，因此从产出的角度看是纯损耗。但痛苦是一种定位信息，它告诉一个人自己的边界在哪里、什么对自己重要、什么状态不可接受。一个人的很多判断，是在经历过某种痛苦之后才获得的，这些判断无法通过阅读关于痛苦的信息来替代，因为阅读给出的是关于痛苦的知识，不是痛苦本身。

疲劳和衰老是另外两个，它们的作用和时间有关。它们把一个人放进时间中，使其状态不可逆。一个系统如果不需要休息，它的状态就不会因为消耗而改变；一个人需要休息，意味着他在某段时间里确实失去了一部分能力，这个损失是真实的。衰老把这个过程拉长到一生，使每个人都知道自己拥有的时间有限。有限性是判断的前提之一：只有时间有限，选择才有意义，一个拥有无限时间的存在不需要在两条路之间做取舍，因此也不需要价值排序。关系和死亡是最后两个，也是最难被纳入任何算法的。一个人的价值在他的关系中被确认，这种确认不是通过信息交换完成的，而是在长期的相互承担中形成的；死亡使一段生命有了整体形状，正是因为它会结束，一生中的选择才有重量。

这些部分不能被折算，不是因为它们太复杂，而是因为它们和产出不在同一个量纲上。把它们折算成效用值，得到的不是它们的等价物，而是它们的一个投影。投影保留了形状，丢掉了深度，而这个丢失不会被察觉，因为投影本身就是完整的图像。

\subhead{负熵这个说法的准确用法}

把人的这部分存在称为负熵体验，需要先交代它的来源，也需要交代它的限度。前面引过的那条物理表述在这里依然成立：生命是一种耗散结构，它靠持续的流动维持一个远离平衡的状态，一旦流动停止，结构就瓦解。\cite{schrodinger1944} 之所以要反复回到这一点，是因为它恰好卡在生成与维持的分界上：生成可以借一次性的投入完成，维持只能靠没有终点的支付。

用这个框架看人，可以得到一个不算新鲜但常被忽略的结论。维持一个活人的低熵状态需要持续的物质和能量投入，而这个投入到死亡为止不会停止。人在这一生中处理的物质和能量总量是可以估算的，与他的产出无关；一个什么也不做的人和一个产出极高的人，在维持自身存在这一项上的消耗在同一个量级。这个框架的价值在于它给出了一个非产出导向的人的定义：人首先是一个需要持续摄取、持续维持、持续老化、最终停止的物理过程。产出只是这个过程的一个分支，而且是可被替代的分支。当分支被替代，主干还在。\key{AI 可以替代许多产出功能，却不能替人承受一个有限生命。}

人的身体与关系需要持续的资源投入，这是处境的一部分。承认这种生活的价值，还需要伦理与制度理由。尊严、平等与参与资格可以支撑相应要求；工作被替代以后，收入、保障与社会承认则需要重新安排。维持成本指出需求所在，权利与分配决定需求怎样被接住。

\subhead{人的价值与评价尺度}

人的价值需要一套能够描述人的尺度。身体、经历、关系、痛苦和选择，都参与一个人的生活；产出核算只记录其中可以交付的部分。技术改变了交付物的生产方式，也就暴露了这套核算的局限。要承认其余部分的价值，需要给出伦理理由与相应制度，并使承认落实为权利、资源和参与。

需要论证的不是机器有没有体验，而是当前把人理解为产出工具的框架不足以描述人。这个论证不依赖任何关于机器的断言，它只依赖对人的描述是否完整。工业化的核算体系把人压缩成产出能力，这个压缩在描述上就是有损的：它丢掉了时间、有限性、关系、承受和历史。丢掉的这些部分在产出导向的社会里可以被忽略，因为它们在核算中为零；当产出被替代，为零的那些部分成了全部剩余，压缩的代价就暴露了。这个论证有一个需要承担的结果：它不保证人一定比机器优越，也不保证人一定有位置。它只说明，如果人只剩产出这一个维度，那么在被替代之后确实什么都不剩；而人实际上不只这一个维度，问题是社会是否有能力承认并核算其余的维度。

前面的讨论给出的是价值会流向哪里，给出的是判断和承受的资格，也给出了个体与组织各自能做的调整。这些结论在描述和实践两个层面都已经落地，它们合起来指向一个更大的问题：这些安排靠什么被固定下来，以及一个文明在决定把资源投向哪里时，依据的是什么。位置与做法都需要一个承载它们的容器，而这个容器不在个体身上，也不在单个组织身上。
